%==============================================================================
% 卒論・修論発表用 Beamer スライドテンプレート - ヘッダーファイル
%
% このファイルにはスライドで使用する各種設定が含まれています。
% slide_main.tex から %==============================================================================
% 卒論・修論発表用 Beamer スライドテンプレート - ヘッダーファイル
%
% このファイルにはスライドで使用する各種設定が含まれています。
% slide_main.tex から %==============================================================================
% 卒論・修論発表用 Beamer スライドテンプレート - ヘッダーファイル
%
% このファイルにはスライドで使用する各種設定が含まれています。
% slide_main.tex から %==============================================================================
% 卒論・修論発表用 Beamer スライドテンプレート - ヘッダーファイル
%
% このファイルにはスライドで使用する各種設定が含まれています。
% slide_main.tex から \input{slide_preamble} で読み込まれます。
% テーマはデフォルト（\usetheme の指定なし）です。
% テーマを変えたい場合は slide_main.tex で \usetheme{...} を指定してください。
%==============================================================================

%------------------ 日本語フォント設定 ----------------------------------------
% 和文をゴシック体（サンセリフ）に統一する。
% Beamer は欧文がサンセリフ体のため、和文もゴシックに合わせると見た目が揃う。
\renewcommand{\kanjifamilydefault}{\gtdefault}

%------------------ ハイパーリンク ---------------------------------------------
% PDF のしおり（ブックマーク）を日本語対応にする（uplatex + dvipdfmx 用）。
% beamer は内部で hyperref を読み込むため、\documentclass の直後に読み込む。
\usepackage{pxjahyper}

%------------------ 数式 --------------------------------------------------------
% beamer は amsmath, amssymb を自動で読み込むため、追加は不要。
% （必要なら以下をコメントアウト解除して追加する）
% \usepackage{mathtools}   % \coloneqq など高度な数式マクロ

%------------------ 画像・表 ----------------------------------------------------
\usepackage{graphicx}    % 画像挿入（\includegraphics）
\usepackage{booktabs}    % 表の罫線（\toprule / \midrule / \bottomrule）
\usepackage{subcaption}  % 複数図を (a), (b) 形式でまとめる（fig:xx-a 等）

%------------------ 動画 --------------------------------------------------------
% 動画を PDF に埋め込む機能は「使用不可」。
% 理由: media9 は Flash Player（終了済み）ベースのため実質再生できないうえ、
%       PDF が約13MBに肥大化する。
% \usepackage{media9} はコメントアウト済み（下の行）。
% どうしても使う場合は slide_main.tex の動画スライド（コメントアウト中）も有効化する。
% \usepackage{media9} % Flash Player死亡のため使用不可
% 使用例:
%   \includemedia[width=0.55\linewidth,height=0.31\linewidth,activate=onclick,
%                 flashvars={autoplay=true}]{ポスター画像}{movie/動画ファイル.mp4}
% 注意: multimedia パッケージの \movie は dvipdfmx と非互換のため使用しないこと

%------------------ 定理環境 ----------------------------------------------------
% beamer には block / alertblock / exampleblock と、
% theorem, lemma, corollary, definition, example, proof が標準で用意されている。
% （スライドでは \begin{block}{タイトル} を主に使うとよい。
%   タイトルは日本語で指定できるので \newtheorem は定義していない。）

%------------------ よく使う数式コマンド ---------------------------------------
\newcommand{\bm}{\boldsymbol}       % 太字ベクトル: \bm{x}
\newcommand{\del}{\partial}         % 偏微分: \del
\newcommand{\diag}{\mathrm{diag}}   % 対角行列: \diag

%------------------ フッター（ページ番号） --------------------------------------
% 右下のナビゲーション記号（ゴミのようなアイコン）を非表示
\setbeamertemplate{navigation symbols}{}
% フッター: 右下に「ページ番号 / 総ページ数」を表示
% （タイトルページや [plain] 指定のフレームには表示されない）
\setbeamertemplate{footline}{%
  \hbox{%
    \begin{beamercolorbox}[wd=\paperwidth,ht=2.5ex,dp=1.125ex]{}%
      \hfill\insertframenumber\,/\,\inserttotalframenumber\hspace{2mm}%
    \end{beamercolorbox}%
  }%
}

%------------------ セクションページ -------------------------------------------
% 各セクションの開始時に、セクションタイトルのスライドを自動挿入する
% 注意: \section*{...} でも発火する（目次には載らないが区切りページは挿入される）
\AtBeginSection{%
  \begin{frame}[plain]
    \sectionpage
  \end{frame}%
}
 で読み込まれます。
% テーマはデフォルト（\usetheme の指定なし）です。
% テーマを変えたい場合は slide_main.tex で \usetheme{...} を指定してください。
%==============================================================================

%------------------ 日本語フォント設定 ----------------------------------------
% 和文をゴシック体（サンセリフ）に統一する。
% Beamer は欧文がサンセリフ体のため、和文もゴシックに合わせると見た目が揃う。
\renewcommand{\kanjifamilydefault}{\gtdefault}

%------------------ ハイパーリンク ---------------------------------------------
% PDF のしおり（ブックマーク）を日本語対応にする（uplatex + dvipdfmx 用）。
% beamer は内部で hyperref を読み込むため、\documentclass の直後に読み込む。
\usepackage{pxjahyper}

%------------------ 数式 --------------------------------------------------------
% beamer は amsmath, amssymb を自動で読み込むため、追加は不要。
% （必要なら以下をコメントアウト解除して追加する）
% \usepackage{mathtools}   % \coloneqq など高度な数式マクロ

%------------------ 画像・表 ----------------------------------------------------
\usepackage{graphicx}    % 画像挿入（\includegraphics）
\usepackage{booktabs}    % 表の罫線（\toprule / \midrule / \bottomrule）
\usepackage{subcaption}  % 複数図を (a), (b) 形式でまとめる（fig:xx-a 等）

%------------------ 動画 --------------------------------------------------------
% 動画を PDF に埋め込む機能は「使用不可」。
% 理由: media9 は Flash Player（終了済み）ベースのため実質再生できないうえ、
%       PDF が約13MBに肥大化する。
% \usepackage{media9} はコメントアウト済み（下の行）。
% どうしても使う場合は slide_main.tex の動画スライド（コメントアウト中）も有効化する。
% \usepackage{media9} % Flash Player死亡のため使用不可
% 使用例:
%   \includemedia[width=0.55\linewidth,height=0.31\linewidth,activate=onclick,
%                 flashvars={autoplay=true}]{ポスター画像}{movie/動画ファイル.mp4}
% 注意: multimedia パッケージの \movie は dvipdfmx と非互換のため使用しないこと

%------------------ 定理環境 ----------------------------------------------------
% beamer には block / alertblock / exampleblock と、
% theorem, lemma, corollary, definition, example, proof が標準で用意されている。
% （スライドでは \begin{block}{タイトル} を主に使うとよい。
%   タイトルは日本語で指定できるので \newtheorem は定義していない。）

%------------------ よく使う数式コマンド ---------------------------------------
\newcommand{\bm}{\boldsymbol}       % 太字ベクトル: \bm{x}
\newcommand{\del}{\partial}         % 偏微分: \del
\newcommand{\diag}{\mathrm{diag}}   % 対角行列: \diag

%------------------ フッター（ページ番号） --------------------------------------
% 右下のナビゲーション記号（ゴミのようなアイコン）を非表示
\setbeamertemplate{navigation symbols}{}
% フッター: 右下に「ページ番号 / 総ページ数」を表示
% （タイトルページや [plain] 指定のフレームには表示されない）
\setbeamertemplate{footline}{%
  \hbox{%
    \begin{beamercolorbox}[wd=\paperwidth,ht=2.5ex,dp=1.125ex]{}%
      \hfill\insertframenumber\,/\,\inserttotalframenumber\hspace{2mm}%
    \end{beamercolorbox}%
  }%
}

%------------------ セクションページ -------------------------------------------
% 各セクションの開始時に、セクションタイトルのスライドを自動挿入する
% 注意: \section*{...} でも発火する（目次には載らないが区切りページは挿入される）
\AtBeginSection{%
  \begin{frame}[plain]
    \sectionpage
  \end{frame}%
}
 で読み込まれます。
% テーマはデフォルト（\usetheme の指定なし）です。
% テーマを変えたい場合は slide_main.tex で \usetheme{...} を指定してください。
%==============================================================================

%------------------ 日本語フォント設定 ----------------------------------------
% 和文をゴシック体（サンセリフ）に統一する。
% Beamer は欧文がサンセリフ体のため、和文もゴシックに合わせると見た目が揃う。
\renewcommand{\kanjifamilydefault}{\gtdefault}

%------------------ ハイパーリンク ---------------------------------------------
% PDF のしおり（ブックマーク）を日本語対応にする（uplatex + dvipdfmx 用）。
% beamer は内部で hyperref を読み込むため、\documentclass の直後に読み込む。
\usepackage{pxjahyper}

%------------------ 数式 --------------------------------------------------------
% beamer は amsmath, amssymb を自動で読み込むため、追加は不要。
% （必要なら以下をコメントアウト解除して追加する）
% \usepackage{mathtools}   % \coloneqq など高度な数式マクロ

%------------------ 画像・表 ----------------------------------------------------
\usepackage{graphicx}    % 画像挿入（\includegraphics）
\usepackage{booktabs}    % 表の罫線（\toprule / \midrule / \bottomrule）
\usepackage{subcaption}  % 複数図を (a), (b) 形式でまとめる（fig:xx-a 等）

%------------------ 動画 --------------------------------------------------------
% 動画を PDF に埋め込む機能は「使用不可」。
% 理由: media9 は Flash Player（終了済み）ベースのため実質再生できないうえ、
%       PDF が約13MBに肥大化する。
% \usepackage{media9} はコメントアウト済み（下の行）。
% どうしても使う場合は slide_main.tex の動画スライド（コメントアウト中）も有効化する。
% \usepackage{media9} % Flash Player死亡のため使用不可
% 使用例:
%   \includemedia[width=0.55\linewidth,height=0.31\linewidth,activate=onclick,
%                 flashvars={autoplay=true}]{ポスター画像}{movie/動画ファイル.mp4}
% 注意: multimedia パッケージの \movie は dvipdfmx と非互換のため使用しないこと

%------------------ 定理環境 ----------------------------------------------------
% beamer には block / alertblock / exampleblock と、
% theorem, lemma, corollary, definition, example, proof が標準で用意されている。
% （スライドでは \begin{block}{タイトル} を主に使うとよい。
%   タイトルは日本語で指定できるので \newtheorem は定義していない。）

%------------------ よく使う数式コマンド ---------------------------------------
\newcommand{\bm}{\boldsymbol}       % 太字ベクトル: \bm{x}
\newcommand{\del}{\partial}         % 偏微分: \del
\newcommand{\diag}{\mathrm{diag}}   % 対角行列: \diag

%------------------ フッター（ページ番号） --------------------------------------
% 右下のナビゲーション記号（ゴミのようなアイコン）を非表示
\setbeamertemplate{navigation symbols}{}
% フッター: 右下に「ページ番号 / 総ページ数」を表示
% （タイトルページや [plain] 指定のフレームには表示されない）
\setbeamertemplate{footline}{%
  \hbox{%
    \begin{beamercolorbox}[wd=\paperwidth,ht=2.5ex,dp=1.125ex]{}%
      \hfill\insertframenumber\,/\,\inserttotalframenumber\hspace{2mm}%
    \end{beamercolorbox}%
  }%
}

%------------------ セクションページ -------------------------------------------
% 各セクションの開始時に、セクションタイトルのスライドを自動挿入する
% 注意: \section*{...} でも発火する（目次には載らないが区切りページは挿入される）
\AtBeginSection{%
  \begin{frame}[plain]
    \sectionpage
  \end{frame}%
}
 で読み込まれます。
% テーマはデフォルト（\usetheme の指定なし）です。
% テーマを変えたい場合は slide_main.tex で \usetheme{...} を指定してください。
%==============================================================================

%------------------ 日本語フォント設定 ----------------------------------------
% 和文をゴシック体（サンセリフ）に統一する。
% Beamer は欧文がサンセリフ体のため、和文もゴシックに合わせると見た目が揃う。
\renewcommand{\kanjifamilydefault}{\gtdefault}

%------------------ ハイパーリンク ---------------------------------------------
% PDF のしおり（ブックマーク）を日本語対応にする（uplatex + dvipdfmx 用）。
% beamer は内部で hyperref を読み込むため、\documentclass の直後に読み込む。
\usepackage{pxjahyper}

%------------------ 数式 --------------------------------------------------------
% beamer は amsmath, amssymb を自動で読み込むため、追加は不要。
% （必要なら以下をコメントアウト解除して追加する）
% \usepackage{mathtools}   % \coloneqq など高度な数式マクロ

%------------------ 画像・表 ----------------------------------------------------
\usepackage{graphicx}    % 画像挿入（\includegraphics）
\usepackage{booktabs}    % 表の罫線（\toprule / \midrule / \bottomrule）
\usepackage{subcaption}  % 複数図を (a), (b) 形式でまとめる（fig:xx-a 等）

%------------------ 動画 --------------------------------------------------------
% 動画を PDF に埋め込む機能は「使用不可」。
% 理由: media9 は Flash Player（終了済み）ベースのため実質再生できないうえ、
%       PDF が約13MBに肥大化する。
% \usepackage{media9} はコメントアウト済み（下の行）。
% どうしても使う場合は slide_main.tex の動画スライド（コメントアウト中）も有効化する。
% \usepackage{media9} % Flash Player死亡のため使用不可
% 使用例:
%   \includemedia[width=0.55\linewidth,height=0.31\linewidth,activate=onclick,
%                 flashvars={autoplay=true}]{ポスター画像}{movie/動画ファイル.mp4}
% 注意: multimedia パッケージの \movie は dvipdfmx と非互換のため使用しないこと

%------------------ 定理環境 ----------------------------------------------------
% beamer には block / alertblock / exampleblock と、
% theorem, lemma, corollary, definition, example, proof が標準で用意されている。
% （スライドでは \begin{block}{タイトル} を主に使うとよい。
%   タイトルは日本語で指定できるので \newtheorem は定義していない。）

%------------------ よく使う数式コマンド ---------------------------------------
\newcommand{\bm}{\boldsymbol}       % 太字ベクトル: \bm{x}
\newcommand{\del}{\partial}         % 偏微分: \del
\newcommand{\diag}{\mathrm{diag}}   % 対角行列: \diag

%------------------ フッター（ページ番号） --------------------------------------
% 右下のナビゲーション記号（ゴミのようなアイコン）を非表示
\setbeamertemplate{navigation symbols}{}
% フッター: 右下に「ページ番号 / 総ページ数」を表示
% （タイトルページや [plain] 指定のフレームには表示されない）
\setbeamertemplate{footline}{%
  \hbox{%
    \begin{beamercolorbox}[wd=\paperwidth,ht=2.5ex,dp=1.125ex]{}%
      \hfill\insertframenumber\,/\,\inserttotalframenumber\hspace{2mm}%
    \end{beamercolorbox}%
  }%
}

%------------------ セクションページ -------------------------------------------
% 各セクションの開始時に、セクションタイトルのスライドを自動挿入する
% 注意: \section*{...} でも発火する（目次には載らないが区切りページは挿入される）
\AtBeginSection{%
  \begin{frame}[plain]
    \sectionpage
  \end{frame}%
}
